\section{Related Work}

\paragraph{Double Descent.}
The non-monotonic relationship between model capacity and generalization was rigorously studied by Belkin et al.~\cite{belkin2019reconciling}, who showed that the classical U-shaped bias-variance curve continues into a second descent in the over-parameterized regime, unifying modern and classical perspectives. Nakkiran et al.~\cite{nakkiran2020deep} provided comprehensive empirical evidence across ResNets, transformers, and random forests, introduced the concept of \emph{effective model complexity}, and demonstrated that the phenomenon is robust across architectures and datasets. Theoretical characterizations have been developed for linear regression~\cite{bartlett2020benign} and kernel methods~\cite{mei2022generalization}, confirming that double descent is not an artifact of neural network non-linearity. Zhang et al.~\cite{zhang2017understanding} challenged classical generalization theory by showing that modern networks can memorize random labels while still generalizing on structured data, setting the stage for understanding the interpolation regime.

\paragraph{Data Augmentation.}
Data augmentation is a standard regularization tool in computer vision. Random cropping and horizontal flipping are near-universal in convolutional network training~\cite{he2016deep,krizhevsky2009learning}. Cutout~\cite{devries2017cutout} randomly masks contiguous image patches, encouraging models to use distributed rather than localized features. MixUp~\cite{zhang2018mixup} linearly interpolates between random pairs of training examples (both inputs and labels), smoothing decision boundaries and flattening the training loss surface. Shorten and Khoshgoftaar~\cite{shorten2019survey} provide a comprehensive survey of augmentation methods. More recent approaches such as CutMix~\cite{yun2019cutmix} and RandAugment~\cite{cubuk2020randaugment} extend these principles further. Despite their widespread use, the mechanistic effect of augmentation on double descent has not been systematically studied.

\paragraph{Loss Landscape and Generalization.}
Li et al.~\cite{li2018visualizing} introduced filter-normalized loss landscape visualization, showing that skip connections in ResNets tend to produce flatter, more navigable loss surfaces. Sharpness-aware minimization~\cite{foret2021sharpness} explicitly seeks flat minima and achieves strong generalization, supporting the empirical link between landscape flatness and test performance. We connect these findings to the double descent setting, hypothesizing that augmentation strategies that flatten the loss landscape also suppress the interpolation peak.

\paragraph{Architecture and Double Descent.}
While double descent has been observed across diverse architectures~\cite{nakkiran2020deep}, the architectural conditions that produce or prevent it are underexplored. Our finding that MLPs exhibit no double descent while ResNets do---under otherwise identical conditions---points to the convolutional inductive bias as a relevant factor, a distinction not examined in prior work.
